% !TeX root = ../papers/nc-p3.tex
\section{二元Hankel空間への非全射な射影と残る二つの捩れ長}
\label{sec:deficient-binary-support}
\begin{theorem}\label{thm:deficient-binary-support}
一文字の直線を落とした一辺の像が $P_5^*\subset P_2^*\otimes P_3^*$ の
真の部分空間 $H$ となる標準正則窓を仮定する。$s=\dim H$ は4か5である。
前節の評価層は一般階数4で、その余核の長さ $t$ は2または3に限られる。
十列の核束は $\mathcal O(-1)^5\oplus\mathcal O(t-7)$、第四文字は $W=M_5$ である。
$H$ を消す各五次式 $f$ について、十列の評価行列 $P$ は二次行列 $A_f$ と
$P A_f=fI_4$ を満たす。五つの一次核列で自明十束を割った束を $J$ とすると、
$H^\perp$ から $\operatorname{Hom}(J,\mathcal O(t-2))$ への単射がある。
特に長さ2の場合、$J$ の自明な直和因子が $6-s$ 個以上必要である。
\end{theorem}
\begin{proof}[圧縮で保持される支持]
落とす文字を $z$ とし、元の辺の射影核を $z\otimes S$ とする。
$\dim S=6-s$ が3以上なら $R_1\cap(S\otimes V^*)\ne0$ であり、その非零元に
$z$ を付けて三次交叉の一次元端支持を得る。従って $s\ge4$。
同じ理由で任意の $s$ について $R_1\cap(S\otimes V^*)=0$ を要する。
圧縮によって $W$ の一文字が消えれば元の対応する三テンソルの $U$ 側の端支持は一次元、
$V$ の一文字が消えれば三テンソルの中央の $U$ 支持が一次元になる。
後者では同じ $z$ を共有する連続純関係ができる。よって $V,W$ とも圧縮後も四次元支持である。
これだけで前節の一般階数3の排除と $\dim M_5\ge4$ の議論が再び使える。

$t=0$ だと $W=M_5$ は一つの八次逆系 $\Phi$ の三次極性空間になり、
$P_5\to M_3$, $f\mapsto\Phi\mathbin{\lrcorner}f$ は同型である。
従って圧縮窓 $\Phi(v)(pqq')$ の $P_5$ 側の階数は6で、真の部分空間ではない。
$1\le t\le3$ では分裂と $W=M_5$ は前節の次数計算で従う。
このとき $f\in H^\perp$ は $fV\subset BP_2$ と同値である。各文字の原像を取って
四本の二次列を並べれば $P A_f=fI_4$ である。
$C$ の零化因子を $Z$ とすると $f$ は当然 $Z$ 上で消える。

核の中の一次直和因子の和 $K_1$ は $K(1)$ の大域評価で内在的に決まる。
$J=\mathcal O^{10}/K_1$ とし、$\mathcal O^{10}$ 上で
$D_f=fI_{10}-A_fP$ を考える。$P D_f=0$ なので像は $K(5)$ に入り、
$K_1$ 上での作用は $fI$ だから商では零である。
従って $D_f$ は $J\to(K/K_1)(5)=\mathcal O(t-2)$ を誘導する。
$K/K_1\hookrightarrow J$ への制限が $f$ の乗法なので、$A_f$ を線型に選んで単射を得る。
$J$ は大域生成だから $t=1$ ではそのような非零写像はない。
$t=2$ の自明因子の下界も束の分裂から従う。
\end{proof}
\begin{proof}[有限長部分の左右同時縮約]
残りの $4\le t\le6$ では $M_5=H^0(C(2))^*$ である。
$C$ の台上の線束を自明化し、零化因子の環を $D$ と書く。
$V$ は $C$ を $D$ 加群として生成し、$W\subset C^*$ は四つの汎関数を与える。
$DW$ が消す $C$ の部分加群を割り、有効な加群 $N$ を取る。
このとき $W$ は $N^*$ を $D$ 加群として生成する。
$N$ の長さは $t$ 以下で、零化因子の次数を $e$ とすれば $e\le6$ である。
五次式はその因子の環へ全射する。圧縮の $P_5$ 側の階数 $s$ は $e$ 以下である。
よって $e\ge4$ となる。

$N$ を各閉点における離散付値環の基本因子に分解する。
各点で最大長の基本因子がすべて二つ以上あれば、長さは最大長の和の二倍、
すなわち $2e$ 以上になり矛盾する。従ってある点 $x$ で最大長 $m$ の因子が一つだけである。
他点で零、$x$ で局所パラメータの $m-1$ 乗となる元を五次式で表す。
この元が $N$ に与える乗法は階数1であり、その像はsocle、その余像は
剰余体商 $N\to k_x$ である。$V$ に制限した余像汎関数 $v_x$ は生成性から非零である。
$W$ の全てがこのsocle上で消えれば $DW$ もそこで消えるので、
その評価による汎関数 $w_x$ も非零である。従ってこの五次式で圧縮窓を縮約すると
$0\ne v_x\otimes w_x\in R_2$ を得る。ここでは圧縮後の $R_2$ の像は元の辺に含まれれば十分である。
一方 $C\twoheadrightarrow N\to k_x$ と $M_3$ への双対包含から
$\operatorname{ev}_x^{(3)}\otimes v_x\in R_1$ である。元の二辺での連続純関係となり排除される。
長さの組合せを\path{computations/verify_deficient_binary_support.py}でも検算する。
\end{proof}

\begin{corollary}\label{cor:no-full-pure-projected-curve}
任意の標準正則な辺で一方の四文字を直線で割ると、射影された関係空間の階数一部分は、
両射影がそれぞれ射影平面と射影三空間を線型に張る整曲線を含まない。
元の辺に純関係があっても同じである。
\end{corollary}
\begin{proof}
その曲線の関数体で両側の座標空間の積は少なくとも $3+4-1=6$ 次元である。
辺は六次元しかないので射影は単射で積次元は最小値になる。
命題\ref{prop:binary-quotient-rank-two-curve}の付値と一次元増大の証明は
この等号だけから両空間を $P_2,P_3$ に同定する。従って前節の定理に反する。
\end{proof}
